\documentclass[10pt,a4paper]{amsart}
\usepackage[margin=19mm]{geometry}
\usepackage{amsmath,amssymb,mathtools}
\usepackage[scr=rsfs,cal=euler]{mathalfa}
\usepackage{xfrac}
\usepackage[unicode,hidelinks,bookmarks=false]{hyperref}
\newcommand{\ie}{i.e.\ }
\newcommand{\cf}{cf.\ }
\newcommand{\resp}{respectively\ }
\newcommand{\Subsection}{\subsection}
\providecommand{\nobreakdash}{\nobreak\hbox{-}}
\setlength{\emergencystretch}{2em}
\multlinegap0pt
\usepackage{amssymb}
\usepackage{xcolor}
\usepackage[shortlabels]{enumitem}
\newtheorem{theorem}{Theorem}
\def \D{\mathbb{D}}
\newcommand{\clb}{\mathcal{B}}
\newcommand{\cle}{\mathcal{E}}
\newcommand{\clf}{\mathcal{F}}
\newcommand{\clm}{\mathcal{M}}
\newcommand{\cln}{\mathcal{N}}
\newcommand{\ran}{\mathrm{ran}}
\title{Odometer maps on Fock spaces: block decompositions, Toeplitz-type realizations, and the adjoint}
\author{Mansi Anil Suryawanshi}
\date{Source: arXiv:2605.26674v1 (CC BY 4.0)}
\begin{document}
\maketitle

\noindent\emph{Source and licence.} Odometer maps on Fock spaces: block decompositions, Toeplitz-type realizations, and the adjoint. Version arXiv:2605.26674v1 (CC BY 4.0), distributed by arXiv under the Creative Commons Attribution 4.0 International licence, http://creativecommons.org/licenses/by/4.0/, as recorded in the arXiv licence metadata for this exact version and shown on the abstract page https://arxiv.org/abs/2605.26674v1. Full licence notice: \texttt{licenses/arxiv-2605.26674-CC-BY-4.0.txt}.\par\smallskip

\noindent\textit{Editorial scope.} Counted unit: the article's theorem thm:operator\_characterization\_symbol (source0.tex lines 1026--1040). The article's complete original proof of it is delivered with it (source0.tex lines 1042--1118, one \texttt{proof} environment of 77 lines). No mathematics is written, repaired or re-derived: the statement and the proof are the article's own lines, in the article's own notation, and no retained line is substituted or re-flowed.  Retained uncounted as the source-local context the proof needs: lines 733--746.  Nothing was omitted from any retained span, so the package carries no omission record.  No retained line contains a citation, so the wrapper carries no bibliography and no bibliography entry.  No retained line contains a figure environment, an image-inclusion command or drawing code, so no image is delivered and the package declares no asset dependency.  Editorial formatting only, in addition: an AMS \texttt{amsart} preamble that restates the article's own declarations and macros used by the retained lines, namely the environments \texttt{theorem} on separate counters, together with the standard packages those lines need.  The retained environments print in this document as Theorem 1 in order of appearance, whereas the article prints the same items with the numbering of its own full document.  The complete native source of the article as distributed in the arXiv e-print for this exact version is bundled unchanged as \texttt{sources/201435/source0.tex}.
\begin{theorem}\label{prop:hardy_space_multiplier}
For \(r\geq 0\), let
\(L_r\in\clb(\cle)\) be defined by $L_r:=(\langle e_1^r,\cdot\rangle\otimes I_\cle)L.$
Then
\[
U_1W_{22}U_n^*=M_\Theta
\quad \text{on } H^2_{\cle}(\D),
\]
where \(M_\Theta\) is the analytic Toeplitz operator with uniquely determined
symbol \(\Theta\in H^\infty_{\clb(\cle)}(\D)\) given by
\[
\Theta(z)=\sum_{r\geq 0}z^rL_r.
\]
\end{theorem}

\begin{theorem}\label{thm:operator_characterization_symbol}
Let $L \in \clb(\cle,\clf_n^2\otimes \cle)$
and let $\Theta \in H^\infty_{\clb(\cle)}(\D)$
be the analytic symbol associated with  $W_L.$
Then 
\begin{enumerate}
    \item \(W_L\) is an isometry if and only if \(L(\cle)\subseteq \clm^\perp\), and \(\Theta\) is inner.

   \item \(W_L\) is unitary if and only if $L\cle\subseteq \clm^\perp$
and \(\Theta\) is a constant unitary operator-valued function.

   \item \(W_L\) is invertible if and only if \(\Theta\) is invertible in
\(H^\infty_{\clb(\cle)}(\D)\). 
 \end{enumerate}
\end{theorem}

\begin{proof}
\begin{enumerate}
        \item 
Suppose that $W_L$ is an isometry. Pick $\eta, \gamma \in \cle$, and let $e_\mu \otimes \gamma \in \clm$. By the definition of the odometer map, there exists a unique preimage $e_{\mu-1} \otimes  \gamma \in \cln$ such that $W_L( e_{\mu-1} \otimes  \gamma)=e_\mu \otimes \gamma$.
Since $W_L$ is an isometry, we have
\[
\langle L\eta,\, e_\mu \otimes \gamma \rangle
=
\langle W_L(\Omega \otimes \eta),\, W_L( e_{\mu-1} \otimes  \gamma) \rangle
= \langle \Omega \otimes \eta,\, e_{\mu-1} \otimes  \gamma \rangle = 0,
\]
where the final equality holds because $e_{\mu-1}\otimes\gamma \in \cln$, and hence is orthogonal to the vacuum space. Hence, $L\cle \subseteq \clm^\perp$, and $W_{12} = 0$. Therefore, $W_L = W_{11} \oplus W_{22}$. Since $W_{11}$ is unitary, it follows that $W_{22}$ is an isometry. As $U_1$ and $U_n$ are unitaries, Theorem~\ref{prop:hardy_space_multiplier} implies
\[
M_{\Theta}^*M_{\Theta}
=
U_n W_{22}^* W_{22} U_n^*
=
I_{H^2_\cle(\D)},
\]
which shows that $M_{\Theta}$ is an isometry. Equivalently, $\Theta$ is inner.

Conversely, if $L(\cle)\subseteq \clm^\perp$ then $W_{12}=0.$ Since $\Theta$ is inner, the Toeplitz operator $M_\Theta$ is an isometry. Therefore, $W_{22}=U_1^*M_\Theta U_n$ is an isometry. Because $W_{11}$ is unitary, the diagonal block operator $W_L = W_{11} \oplus W_{22}$ is an isometry.

\item By (1) above \(W_L\) is unitary if
and only if $L\cle\subseteq\clm^\perp$ and \(M_\Theta\) is unitary, which holds if and only if \(\Theta\) is
a constant unitary operator-valued function.
\item
Since $M_\Theta = U_1 W_{22} U_n^*,$ we have $W_{22}\text{ is invertible}$ if and only if $
M_\Theta\text{ is invertible}.$
Equivalently, \(\Theta\) is invertible in
\(H^\infty_{\clb(\cle)}(\D)\), that is,
\(\Theta^{-1}\in H^\infty_{\clb(\cle)}(\D)\).
Assume first that $\Theta \in H^\infty_{\clb(\cle)}(\D)$ is invertible. Define $X:\clm\oplus \clm^\perp \to \cln\oplus \cln^\perp$ by
\[
X:=
\begin{pmatrix} 
W_{11}^* & -\,W_{11}^*W_{12}W_{22}^{-1}\\
0 & W_{22}^{-1}
\end{pmatrix}.
\]
Since $W_{11}$ is unitary, a direct computation gives
\[
W_LX=
\begin{pmatrix}
W_{11} & W_{12}\\
0 & W_{22}
\end{pmatrix}
\begin{pmatrix}
W_{11}^* & -\,W_{11}^*W_{12}W_{22}^{-1}\\
0 & W_{22}^{-1}
\end{pmatrix}
=
\begin{pmatrix}
I_{\clm} & 0\\
0 & I_{\clm^\perp}
\end{pmatrix}.
\]
That is, $W_LX=I_{\clm\oplus \clm^\perp}$. A similar computation implies $XW_L= I_{\cln\oplus \cln^\perp}.$
Thus $X=W_L^{-1}$, and hence $W_L$ is invertible.


Conversely, assume that $W_L$ is invertible. We prove that $W_{22}:\cln^{\perp} \to \clm^{\perp}$ is both surjective and injective.
First, let $y\in \clm^\perp$. Since $W_L$ is onto, there exists $u\oplus v\in \cln\oplus \cln^\perp$ such that $W_L(u\oplus v)=0\oplus y.$
Equivalently,
\[
W_{11}u+W_{12}v=0, \qquad W_{22}v=y.
\]
Hence $y\in \ran W_{22}$, so $W_{22}$ is surjective.
Next, let $v\in \cln^\perp$ satisfy $W_{22}v=0$. Since $W_{11}:\cln\to \clm$ is unitary, there exists $u\in \cln$ such that $W_{11}u=-\,W_{12}v.$
Then
\[
W_L(u\oplus v)=\big(W_{11}u+W_{12}v\big)\oplus W_{22}v=0\oplus 0.
\]
As $W_L$ is injective, it follows that $u=0$ and $v=0$. Hence $\ker W_{22}=\{0\}$. Since $W_{22}: \cln^\perp \to \clm^\perp$ is a bounded bijection between Hilbert spaces, the bounded inverse theorem implies that \(W_{22}^{-1}\in \clb(\clm^\perp,\cln^\perp)\).
Consequently $\Theta$ is invertible. 
\end{enumerate}
\end{proof}

\end{document}
